\documentclass{amsart}
% For dealing with references we use the comment environment
\usepackage{verbatim}
\newenvironment{reference}{\comment}{\endcomment}
%\newenvironment{reference}{}{}
\newenvironment{slogan}{\comment}{\endcomment}
\newenvironment{history}{\comment}{\endcomment}

% For commutative diagrams we use Xy-pic
\usepackage[all]{xy}

% We use 2cell for 2-commutative diagrams.
\xyoption{2cell}
\UseAllTwocells

% We use multicol for the list of chapters between chapters
\usepackage{multicol}

% This is generally recommended for better output
\usepackage{lmodern}
\usepackage[T1]{fontenc}

% For cross-file-references

% Package for hypertext links:
\usepackage{hyperref}

% For any local file, say "hello.tex" you want to link to please

% Theorem environments.
%
\theoremstyle{plain}
\newtheorem{theorem}[subsection]{Theorem}
\newtheorem{proposition}[subsection]{Proposition}
\newtheorem{lemma}[subsection]{Lemma}

\theoremstyle{definition}
\newtheorem{definition}[subsection]{Definition}
\newtheorem{example}[subsection]{Example}
\newtheorem{exercise}[subsection]{Exercise}
\newtheorem{situation}[subsection]{Situation}

\theoremstyle{remark}
\newtheorem{remark}[subsection]{Remark}
\newtheorem{remarks}[subsection]{Remarks}

\numberwithin{equation}{subsection}

% Macros
%
\def\lim{\mathop{\mathrm{lim}}\nolimits}
\def\colim{\mathop{\mathrm{colim}}\nolimits}
\def\Spec{\mathop{\mathrm{Spec}}}
\def\Hom{\mathop{\mathrm{Hom}}\nolimits}
\def\Ext{\mathop{\mathrm{Ext}}\nolimits}
\def\SheafHom{\mathop{\mathcal{H}\!\mathit{om}}\nolimits}
\def\SheafExt{\mathop{\mathcal{E}\!\mathit{xt}}\nolimits}
\def\Sch{\mathit{Sch}}
\def\Mor{\mathop{\mathrm{Mor}}\nolimits}
\def\Ob{\mathop{\mathrm{Ob}}\nolimits}
\def\Sh{\mathop{\mathit{Sh}}\nolimits}
\def\NL{\mathop{N\!L}\nolimits}
\def\CH{\mathop{\mathrm{CH}}\nolimits}
\def\proetale{{pro\text{-}\acute{e}tale}}
\def\etale{{\acute{e}tale}}
\def\QCoh{\mathit{QCoh}}
\def\Ker{\mathop{\mathrm{Ker}}}
\def\Im{\mathop{\mathrm{Im}}}
\def\Coker{\mathop{\mathrm{Coker}}}
\def\Coim{\mathop{\mathrm{Coim}}}

% Boxtimes
%
\DeclareMathSymbol{\boxtimes}{\mathbin}{AMSa}{"02}

%
% Macros for moduli stacks/spaces
%
\def\QCohstack{\mathcal{QC}\!\mathit{oh}}
\def\Cohstack{\mathcal{C}\!\mathit{oh}}
\def\Spacesstack{\mathcal{S}\!\mathit{paces}}
\def\Quotfunctor{\mathrm{Quot}}
\def\Hilbfunctor{\mathrm{Hilb}}
\def\Curvesstack{\mathcal{C}\!\mathit{urves}}
\def\Polarizedstack{\mathcal{P}\!\mathit{olarized}}
\def\Complexesstack{\mathcal{C}\!\mathit{omplexes}}
% \Pic is the operator that assigns to X its picard group, usage \Pic(X)
% \Picardstack_{X/B} denotes the Picard stack of X over B
% \Picardfunctor_{X/B} denotes the Picard functor of X over B
\def\Pic{\mathop{\mathrm{Pic}}\nolimits}
\def\Picardstack{\mathcal{P}\!\mathit{ic}}
\def\Picardfunctor{\mathrm{Pic}}
\def\Deformationcategory{\mathcal{D}\!\mathit{ef}}

\setlength{\emergencystretch}{3em}
\begin{document}
\title[Stacks Project, Tag 0AU9]{260 --- Existence of normalized dualizing complexes on finite-type schemes}
\maketitle
\noindent Copyright (C) 2005--2025 Johan de Jong. Stacks Project authors. Source: \href{https://stacks.math.columbia.edu/tag/0AU9}{Tag 0AU9}.

\noindent Editorial difficulty note (not author text). Difficulty H2: The proof inductively glues derived dualizing objects over a finite separated covering. It uses normalized uniqueness to obtain overlap identifications and then verifies compatibility of the new isomorphisms with every transition map. The substantive obligation is an effective derived gluing construction with normalization..

\noindent Editorial provenance note (not author text). History: excerpt from revision \texttt{a0444\allowbreak 6e57e\allowbreak c1fbc\allowbreak 252a8\allowbreak 71afc\allowbreak ec775\allowbreak 2fb28\allowbreak 07b14}, duality.tex, lines 5174--5229. Reference presentation was adapted; original-author context and core support blocks were retained or added. The mathematical wording is unchanged. Licensed under GNU FDL 1.2 or later; no invariant sections or cover texts. See ../licenses/COPYING. Context blocks reproduce author definitions and supporting results; they are not additional counted problems.

\begin{situation}
\label{situation-dualizing}
Here $S$ is a Noetherian scheme and $\omega_S^\bullet$ is a dualizing
complex.
\end{situation}

\noindent
In Situation \ref{situation-dualizing} let $X$ be a scheme of finite type
over $S$. Let $\mathcal{U} : X = \bigcup_{i = 1, \ldots, n} U_i$
be a finite open covering of $X$ by objects of $\textit{FTS}_S$, see
Situation \ref{situation-shriek}. All this means is that the morphisms
$U_i \to S$ are separated (as they are already of finite type).
Every affine scheme of finite type over $S$ is an object of $\textit{FTS}_S$
by Schemes, Lemma \href{https://stacks.math.columbia.edu/tag/01KV}{lemma compose after separated (Tag 01KV)}
hence such open coverings certainly exist.
Then for each $i, j, k \in \{1, \ldots, n\}$
the morphisms $p_i : U_i \to S$, $p_{ij} : U_i \cap U_j \to S$,
and $p_{ijk} : U_i \cap U_j \cap U_k \to S$ are separated and
each of these schemes is an object of $\textit{FTS}_S$.
From such an open covering we obtain
\begin{enumerate}
\item $\omega_i^\bullet = p_i^!\omega_S^\bullet$
a dualizing complex on $U_i$, see Section \href{https://stacks.math.columbia.edu/tag/0AU3}{section duality (Tag 0AU3)},
\item for each $i, j$ a canonical isomorphism
$\varphi_{ij} :
\omega_i^\bullet|_{U_i \cap U_j} \to \omega_j^\bullet|_{U_i \cap U_j}$, and
\item
\label{item-cocycle-glueing}
for each $i, j, k$ we have
$$
\varphi_{ik}|_{U_i \cap U_j \cap U_k} =
\varphi_{jk}|_{U_i \cap U_j \cap U_k} \circ
\varphi_{ij}|_{U_i \cap U_j \cap U_k}
$$
in $D(\mathcal{O}_{U_i \cap U_j \cap U_k})$.
\end{enumerate}
Here, in (2) we use that $(U_i \cap U_j \to U_i)^!$
is given by restriction (Lemma \href{https://stacks.math.columbia.edu/tag/0AU0}{lemma shriek open immersion (Tag 0AU0)})
and that we have canonical isomorphisms
$$
(U_i \cap U_j \to U_i)^! \circ p_i^! = p_{ij}^! =
(U_i \cap U_j \to U_j)^! \circ p_j^!
$$
by Lemma \href{https://stacks.math.columbia.edu/tag/0ATX}{lemma upper shriek composition (Tag 0ATX)} and to get (3) we use
that the upper shriek functors form a pseudo functor by
Lemma \href{https://stacks.math.columbia.edu/tag/0ATY}{lemma pseudo functor (Tag 0ATY)}.

\medskip\noindent
In the situation just described a
{\it dualizing complex normalized relative to $\omega_S^\bullet$
and $\mathcal{U}$} is a pair $(K, \alpha_i)$ where $K \in D(\mathcal{O}_X)$
and $\alpha_i : K|_{U_i} \to \omega_i^\bullet$ are isomorphisms
such that $\varphi_{ij}$ is given by
$\alpha_j|_{U_i \cap U_j} \circ \alpha_i^{-1}|_{U_i \cap U_j}$.
Since being a dualizing complex on a scheme is a local property
we see that dualizing complexes normalized relative to $\omega_S^\bullet$
and $\mathcal{U}$ are indeed dualizing complexes.



\begin{situation}
\label{situation-shriek}
Here $S$ is a Noetherian scheme and $\textit{FTS}_S$ is the category whose
\begin{enumerate}
\item objects are schemes $X$ over $S$ such that the structure
morphism $X \to S$ is both separated and of finite type, and
\item morphisms $f : X \to Y$ between objects are morphisms
of schemes over $S$.
\end{enumerate}
\end{situation}

\begin{lemma}
\label{lemma-existence-good-dualizing}
In Situation \ref{situation-dualizing} let $X$ be a scheme of finite type
over $S$ and let $\mathcal{U}$ be a finite open covering
of $X$ by schemes separated over $S$. Then there exists
a dualizing complex normalized relative to $\omega_S^\bullet$ and
$\mathcal{U}$.
\end{lemma}

\begin{proof}
Say $\mathcal{U} : X = \bigcup_{i = 1, \ldots, n} U_i$.
We prove the lemma by induction on $n$. The base case $n = 1$ is immediate.
Assume $n > 1$. Set $X' = U_1 \cup \ldots \cup U_{n - 1}$
and let $(K', \{\alpha'_i\}_{i = 1, \ldots, n - 1})$
be a dualizing complex normalized relative to $\omega_S^\bullet$
and $\mathcal{U}' : X' = \bigcup_{i = 1, \ldots, n - 1} U_i$.
It is clear that $(K'|_{X' \cap U_n}, \alpha'_i|_{U_i \cap U_n})$
is a dualizing complex normalized relative to $\omega_S^\bullet$
and the covering
$X' \cap U_n = \bigcup_{i = 1, \ldots, n - 1} U_i \cap U_n$.
On the other hand, by condition (\ref{item-cocycle-glueing}) the pair
$(\omega_n^\bullet|_{X' \cap U_n}, \varphi_{ni})$
is another dualizing complex normalized relative to $\omega_S^\bullet$
and the covering
$X' \cap U_n = \bigcup_{i = 1, \ldots, n - 1} U_i \cap U_n$.
By Lemma \href{https://stacks.math.columbia.edu/tag/0AU7}{lemma good dualizing unique (Tag 0AU7)} we obtain a unique isomorphism
$$
\epsilon : K'|_{X' \cap U_n} \longrightarrow \omega_i^\bullet|_{X' \cap U_n}
$$
compatible with the given local isomorphisms.
By Cohomology, Lemma \href{https://stacks.math.columbia.edu/tag/08DG}{lemma glue (Tag 08DG)}
we obtain $K \in D(\mathcal{O}_X)$ together with
isomorphisms $\beta : K|_{X'} \to K'$ and
$\gamma : K|_{U_n} \to \omega_n^\bullet$ such that
$\epsilon = \gamma|_{X'\cap U_n} \circ \beta|_{X' \cap U_n}^{-1}$.
Then we define
$$
\alpha_i = \alpha'_i \circ \beta|_{U_i}, i = 1, \ldots, n - 1,
\text{ and }
\alpha_n = \gamma
$$
We still need to verify that $\varphi_{ij}$ is given by
$\alpha_j|_{U_i \cap U_j} \circ \alpha_i^{-1}|_{U_i \cap U_j}$.
For $i, j \leq n - 1$ this follows from the corresponding
condition for $\alpha_i'$. For $i = j = n$ it is clear as well.
If $i < j = n$, then we get
$$
\alpha_n|_{U_i \cap U_n} \circ \alpha_i^{-1}|_{U_i \cap U_n} =
\gamma|_{U_i \cap U_n} \circ \beta^{-1}|_{U_i \cap U_n}
\circ (\alpha'_i)^{-1}|_{U_i \cap U_n} =
\epsilon|_{U_i \cap U_n} \circ (\alpha'_i)^{-1}|_{U_i \cap U_n}
$$
This is equal to $\alpha_{in}$ exactly because $\epsilon$
is the unique map compatible with the maps
$\alpha_i'$ and $\alpha_{ni}$.
\end{proof}
\end{document}
